\chapter{绪论}\label{chap:1}

\section{背景介绍}\label{1.1}

近年来，神经退行性疾病与精神疾病的发病率和社会负担持续上升，推动了脑科学与神经医学领域对脑功能机制及相关病理过程的深入研究\cite{RN2016}。在中国，人口老龄化进程不断加快，同时工业化和城镇化所带来的环境与生活方式变化进一步增加了帕金森病(Parkinson's Disease, PD)等神经系统疾病的防治压力\cite{RN2015}。因此，围绕神经活动调控、能量代谢异常以及疾病进展机制开展高精度研究，已成为脑科学研究中的重要课题。

这类研究的核心需求在于：能够在活体条件下，对深层组织中的神经活动与代谢过程进行连续、动态且可靠的观测。由于脑组织具有高度散射性和复杂的空间结构，传统检测手段往往难以同时兼顾成像深度、时空分辨率以及定量分析能力。特别是在神经元活动、血流代谢耦合以及细胞微环境变化等研究中，获取具有高对比度、高特异性和高稳定性的成像结果，对于揭示脑功能运行规律和疾病发生机制具有重要意义。

因此，发展兼具较大成像深度、高时空分辨率以及定量表征能力的活体显微成像技术，不仅是基础神经科学研究的重要技术支撑，也是推动脑疾病早期诊断、疗效评估和机制研究的关键途径之一。
\section{活体成像探测手段简介}\label{1.2}

面向活体脑成像研究，当前已经发展出多种探测与成像方法。总体上，这些技术可分为非光学探测手段和光学成像手段两大类。前者通常具有较大的成像深度、较强的整体尺度表征能力，适合脑结构、脑功能网络及代谢活动的宏观观测；后者则在空间分辨率、分子特异性以及动态过程记录方面具有明显优势，更适用于细胞尺度乃至亚细胞尺度的神经活动研究。不同技术路线具有不同的适用范围和性能边界，共同构成了现代脑科学研究的重要技术基础。

由于本文的研究重点主要围绕高分辨率活体光学显微成像及其相关探测方法展开，因此本节不对所有脑成像技术进行全面综述，而是结合后续研究主线，对与本文密切相关的几类代表性成像手段进行归纳介绍，并在此基础上引出本文所采用的技术平台与研究意义。

\subsection{非光学探测手段}\label{1.2.1}

非光学探测技术为脑科学研究提供了从整体到局部、从结构到功能的多尺度观测能力，典型代表包括磁共振成像(Magnetic Resonance Imaging, MRI)\cite{RN2107}、功能性磁共振成像(Functional Magnetic Resonance Imaging, fMRI)\cite{RN2006}、正电子发射断层扫描(Positron Emission Tomography, PET)\cite{RN2008}、单光子发射计算机断层扫描(Single-Photon Emission Computed Tomography, SPECT)\cite{RN2009}以及电生理技术\cite{RN2011}等。

其中，MRI 与 fMRI 具有非侵入、成像深度大以及可实现全脑范围观测等优点，在脑结构成像和脑功能连接研究中得到广泛应用。然而，这类方法通常难以兼顾细胞尺度分辨率与分子特异性，因而不适于直接表征微观尺度下的神经元活动过程。PET 与 SPECT 依赖放射性示踪剂获取代谢或分子层面的功能信息，在疾病诊断与示踪研究中具有独特优势，但其空间分辨率相对有限，且实验过程对示踪剂与成像设备依赖较强。电生理方法，如膜片钳和多通道电极记录，则具有极高的时间分辨率，能够直接获取神经元电活动信号，但此类方法往往侵入性较强，观测范围有限，也不利于实现长期、稳定的大范围活体成像。

总体而言，非光学探测手段在脑科学研究中具有不可替代的重要价值，但在高空间分辨率、深层动态观测以及分子特异性表征等方面仍存在一定局限。这些限制进一步推动了高分辨率活体光学成像技术的发展，使其成为揭示脑组织微观结构与功能活动的重要手段。

\subsection{光学成像手段}\label{1.2.2}

\begin{figure}[htbp]
    \centering
    \includegraphics[width=0.75\textwidth]{Fig/Fig01.png}
    \bicaption{用于细胞尺度神经观测的代表性光学成像技术原理对比示意图}{Comparative schematic of representative optical imaging modalities for cellular-scale neural observation}
    \label{fig01}
\end{figure}

相比非光学探测手段，光学成像能够以更高的空间分辨率和更强的分子特异性记录神经活动，是连接细胞机制与网络功能的重要技术路径。当前，应用于活体神经成像的代表性光学方法包括宽场荧光成像(Wide-field Imaging)\cite{RN2017,RN131,RN1805}、光片显微成像(Light-sheet Microscopy, LSM)\cite{RN2113}、荧光光度计(In-vivo Fluorescence Photometry)\cite{RN2114}、光学相干层析成像(Optical Coherence Tomography, OCT)\cite{RN2108}、光声成像(Photoacoustic Imaging, PAI)\cite{RN1464}、共聚焦荧光显微成像(Confocal Fluorescence Microscopy, CFM)\cite{RN1326,RN1938}以及多光子显微成像(Multi-Photon Microscopy, MPM)\cite{RN2024,RN2028}等。

如图\ref{fig01}所示，针对细胞尺度神经观测，已经发展出多种具有代表性的光学成像技术。宽场荧光成像采用整体照明并通过相机直接记录整个视场的荧光信号，具有成像速度快、系统结构简单等优点，但其离焦背景较强，通常需要结合去卷积算法提升图像对比度与轴向分辨率。共聚焦显微成像采用点扫描激发，并利用针孔抑制焦平面外的离焦和散射信号，从而实现较好的光学切片能力和空间分辨率。多光子显微成像利用近红外超快脉冲激光在焦点处产生非线性局域激发，仅在焦点附近形成有效荧光信号，因此具有内禀光学切片能力，同时在深层组织中表现出更好的穿透能力和更低的离焦光损伤。

除上述方法外，基于干涉的结构光照明显微成像(Structured Illumination Microscopy, SIM)通过条纹照明和频谱重建突破衍射极限，可获得超分辨图像；光片显微成像利用与探测光路正交的薄光片照明实现低光毒性和快速体成像；晶格光片显微成像(Lattice Light-sheet, LLS)则在光片显微基础上进一步提升了光片质量与轴向分辨率。另一方面，单分子定位显微成像(Single-Molecule Localization Microscopy, SMLM)通过稀疏发射与高精度定位实现超分辨重建，受激发射损耗显微成像(Stimulated Emission Depletion, STED)则通过环形损耗光束压缩有效发光体积，从而突破衍射极限。这些技术在超分辨和快速三维观测方面各具优势，但在活体深层脑组织中常常受到散射增强、背景升高、成像深度有限、重建复杂度高或系统实现难度大等因素的限制。

对于活体脑组织尤其是原生深层组织成像而言，成像系统不仅需要具有较高的空间分辨率和光学切片能力，还需要在较强散射环境下维持较好的信号保真度与较低的光毒性。在这一背景下，共聚焦显微镜和多光子显微镜成为细胞尺度活体脑成像中最具代表性的两类技术路线。

共聚焦显微镜通过针孔结构有效抑制离焦背景，具有较高的横向分辨率和良好的层析能力，在浅层组织成像和高对比度成像中具有明显优势。然而，针孔在滤除离焦光的同时，也会阻挡部分虽经散射但仍携带有效空间信息的光子，从而造成信号损失。随着成像深度增加，组织散射、折射率不匹配(Refractive Index Mismatch)以及球差(Spherical Aberration)会共同导致点扩散函数(Point Spread Function, PSF)展宽、焦点质量下降，进一步限制其在深层组织中的成像能力\cite{RN2022}。若通过提高激发功率或延长曝光时间补偿信号衰减，则又容易引入更严重的光漂白和光毒性问题。

相比之下，多光子显微成像采用近红外飞秒脉冲激发，依赖非线性局域激发机制仅在焦点附近产生有效荧光信号，因此无需针孔即可获得内禀光学切片能力。同时，由于近红外光在生物组织中具有更低的散射和更深的穿透能力，多光子显微成像在深层组织观测中能够保持更好的信号保真度，并显著降低离焦区域的光损伤\cite{RN2024,RN2028}。因此，共聚焦与多光子显微成像构成了当前活体细胞尺度神经观测的核心技术平台，也是本文开展荧光寿命成像研究的主要成像基础。

\subsection{活体成像中的荧光探针}\label{1.2.3}

\begin{figure}[htbp]
    \centering
    \includegraphics[width=0.9\textwidth]{Fig/Fig02.jpg}
    \bicaption{不同荧光分子对应的激发波长范围示意图}{Schematic of excitation wavelength ranges for different fluorescent molecules}
    \label{fig02}
\end{figure}

在活体神经成像研究中，荧光探针是连接生物过程与光学信号的重要媒介。其中，钙离子作为神经元活动最重要的胞内信号之一，在静息状态下通常维持较低浓度，而在动作电位发放或突触活动发生时，其胞内浓度会迅速升高。钙成像正是利用这一生理过程，通过荧光探针将与神经活动相关的钙信号转化为可探测的光学信号，从而实现对神经活动的间接观测。

目前，神经成像实验中常用的荧光探针主要包括化学钙染料和基因编码荧光指示剂\cite{RN256}。不同探针在光谱特性、标记方式、生物相容性以及适用场景等方面均存在差异，因此对应的激发条件和成像性能也不尽相同。对于多光子显微成像系统而言，探针的选择不仅影响荧光信号强度和可实现的成像深度，还会直接影响激发波长范围、激光脉冲参数以及整个系统光路架构的设计。

如图\ref{fig02}所示，不同类型的荧光分子，包括常见荧光染料、荧光蛋白、钙指示剂以及部分光遗传相关分子，在多光子激发条件下通常对应不同的最佳激发波长范围。与此同时，不同飞秒激光器的可调谐范围和输出特性并不相同，因此在具体系统设计中，需要根据目标探针的光谱响应特性，合理匹配激发波段与激光器参数。对于面向神经活动记录的多光子成像系统而言，这种“探针—激发波长—激光器”之间的协同匹配关系，是确定成像方案、优化系统性能并实现高质量成像的重要依据。

进一步而言，探针的选择还会影响后续寿命成像和功能成像的定量能力。不同荧光分子的荧光寿命、量子效率、光稳定性及环境敏感性存在显著差异，这些因素不仅决定了信号采集效率，也会影响寿命解算的准确性和最终图像质量。因此，在构建面向活体脑成像的荧光寿命显微系统时，探针特性并非独立于光学系统之外的附属条件，而是需要与激发方式、探测链路以及信号处理算法共同考虑的重要组成部分。



\section{文章组织结构}\label{1.3}
本文围绕``共聚焦与多光子荧光寿命成像系统构建''展开研究，全文共分为六章，各章内容安排如下。

第二章阐述显微成像系统设计的理论基础，主要包括点扫描显微镜的基本原理、扫描结构以及系统视场设计方法。

第三章介绍激光扫描共聚焦显微镜系统的构建，主要包括共聚焦成像原理以及多套多色共聚焦荧光寿命成像系统的设计与实现。

第四章介绍多光子显微镜系统的构建，主要包括多光子显微成像的基本原理、两套多光子显微镜系统的搭建，以及多光子显微系统中的脉冲压缩设计。

第五章介绍荧光寿命成像系统的实现，主要包括荧光寿命成像技术概述、系统实现细节，以及基于两种不同方法的荧光寿命获取与结果分析。

第六章介绍本文开展的相关辅助工作，主要包括一套用于超快脉冲表征的自相关仪搭建，以及一套基于双椭圆结构的扩大扫描视场设计。